\scanpage{33}
a) $\vec A\hookrightarrow S\times A$, \deleted{$\unsure{A|\forall}$} $\vec A\to A$ de rang $2$

a') $A^\uparrow\hookrightarrow A\times F$, \quad $A^\uparrow\to A$ de rang $2$

\diagram{\begin{tikzcd}[ampersand replacement=\&,column sep=tiny,row sep=small]
\&\vec A\arrow[dl]\arrow[dr]\&\&A^\uparrow\arrow[dl]\arrow[dr]\&\\
S\arrow[d,equal]\&\& A\arrow[d,equal]\&\&F\arrow[d,equal]\\
\mathcal X_0\&\&\mathcal X_1\&\&\mathcal X_2
\end{tikzcd}}

b) $A^{\vcenter{\hbox{\tikz[x=.15cm,y=.15cm]{\draw[->] (0,0)--(0,1.2);\draw[->] (0,0)--(1.3,0);}}}}=\vec A\times_A A^\uparrow\to S\times F$ est de rang $2$ sur son image

c) $\forall s\in S$, l'ens $\mathcal X_{\inserted{>s}}\deleted{\unsure{(s)}}$ des $a\in A$ et $f\in F$
incidents à $s$ (qui est un \uline{graphe} de
sommets $A$, d'arêtes $F$, grâce à b), et
\unsure{$=$|$\subset$|$\simeq$} un \uline{contour} grâce à a')) est un
\uline{polygône} (i.e. est un contour \inserted{$\ne\varnothing$} \uline{connexe})

c') $\forall f\in F$, l'ens $\mathcal X_{<f}$ des $a\in A$ et $s\in S$
incidents à $f$ (qui est un \uline{graphe} \ldots{}
grâce à a), et un \uline{contour} grâce à b))
est un \uline{polygône} (i.e. est un contour \inserted{$\ne\varnothing$} \uline{conn.})

Version commune (introduisant $\mathcal X_{-1}=\{\varnothing\}$, $\mathcal X_3=\{1\}$ (plus \unsure{grds|grand}
et plus petit élément)

Si $-1\leq i\unsure{<|\leq}j\leq3$, \inserted{avec} $\delta=j-i\geq2$, et \deleted{\unsure{et}} $(i,j)\ne(-1,3)$, si $x_i\in\mathcal X_i$, $y_j\in\mathcal X_j$,
et $\mathcal X_{i<\,?\,<j}$ est l'ens ordonné des $x\in\mathcal X$ \deleted{$\unsure{<}$} tels que
$x_i<x<y_j$, alors la réalisation topologique de $\mathcal X_{i\leq\,\,\leq j}$ est
homéomorphe à \deleted{$\unsure{\Sigma}$} $\unsure{S|\mathbb S}^{\delta-1}$ (sphère de dim $\delta-1$). [Se
généralise aux polyèdres lisses de dim quelconque, cf plus
bas\ldots]

Les conditions a), a'), b) permettent de définir sur
$A^{\vcenter{\hbox{\tikz[x=.15cm,y=.15cm]{\draw[->] (0,0)--(0,1.2);\draw[->] (0,0)--(1.3,0);}}}}$ les opérations $\sigma_1,\sigma_2$ et $\varepsilon$ [avec $\sigma_1\sigma_2=\sigma_2\sigma_1$ \deleted{\unsure{et}}, soit $\sigma$,
et $\sigma_1,\sigma_2,\varepsilon$ étant des involutions, $\sigma_1,\sigma_2,\sigma,\varepsilon$ étant ss
pts fixes ---] les conditions c), c') assurent que
$S,F$ ``sont'' bien les ens de sommets et des faces de la
carte lisse combinatoire ainsi définie\ldots
% archive sheet number: 32
